\section{Model Construction}

\subsection{Protection Definition}

We define protection as multi-dimensional coverage that integrates spatial, temporal, and threat-based criteria. A location $i$ is considered protected at time $t$ if it satisfies:

\begin{equation}
P_i(t) = \mathbb{I}\left[C_{spatial}(i,t) \land C_{temporal}(i,t) \land C_{threat}(i,t)\right] \geq \tau
\label{eq:protection}
\end{equation}

where:
\begin{itemize}
    \item $C_{spatial}(i,t)$ represents spatial coverage within effective monitoring range
    \item $C_{temporal}(i,t)$ ensures temporal coverage during relevant activity periods
    \item $C_{threat}(i,t)$ addresses threat-specific response capabilities
    \item $\tau$ is the protection threshold (typically $\tau = 0.7$)
\end{itemize}

\subsection{Priority Framework}

Species priority combines three critical dimensions: IUCN conservation status, ecological role, and flagship value. We formulate the priority score $w_s$ for species $s$ as:

\begin{equation}
w_s = \alpha \cdot \text{IUCN}_s + \beta \cdot \text{Ecological}_s + \gamma \cdot \text{Flagship}_s
\label{eq:priority}
\end{equation}

where $\alpha + \beta + \gamma = 1$ with empirical values $\alpha = 0.5$, $\beta = 0.3$, $\gamma = 0.2$.

High-priority species identified through this framework include:
\begin{itemize}
    \item \textbf{Black rhino} (\textit{Diceros bicornis}): Critically Endangered, keystone grazer
    \item \textbf{Elephant} (\textit{Loxodonta africana}): Endangered, ecosystem engineer
    \item \textbf{Lion} (\textit{Panthera leo}): Vulnerable, apex predator, high flagship value
\end{itemize}

\subsection{Optimization Problem}

Our core optimization maximizes system-level protection efficiency subject to resource constraints:

\begin{equation}
\max_{x \in \mathcal{X}} \quad \mathcal{P}_{system} = \frac{\sum_{i} w_i \cdot P_i(x)}{\sum_{i} w_i}
\label{eq:optimization}
\end{equation}

subject to:
\begin{align}
\sum_{j} c_j x_j &\leq B_{total} \quad \text{(Budget constraint)} \\
\sum_{k} r_k &\leq R_{available} \quad \text{(Personnel constraint)} \\
P_i(x) &\geq \tau_{min} \quad \forall i \in \mathcal{H} \quad \text{(Minimum protection for high-priority zones)}
\end{align}

where:
\begin{itemize}
    \item $x \in \mathcal{X}$ represents the allocation vector across monitoring methods
    \item $w_i$ is the priority weight for location $i$
    \item $c_j$ are unit costs for monitoring method $j$
    \item $r_k$ represents ranger requirements
    \item $\mathcal{H}$ denotes the set of high-priority zones
\end{itemize}

\subsection{Monitoring Methods}

We consider three primary monitoring methods, each with distinct capabilities:

\begin{table}[h]
\centering
\caption{Monitoring Method Characteristics}
\begin{tabular}{lccc}
\hline
\textbf{Method} & \textbf{Spatial Range} & \textbf{Temporal Coverage} & \textbf{Cost Index} \\
\hline
Ranger patrols & 5-10 km/day/shift & High (8-12 hours) & 1.0 (baseline) \\
UAV surveillance & 20-50 km/flight & Medium (2-4 hours) & 0.3 \\
Camera traps & Point location & Low (event-triggered) & 0.1 \\
\hline
\end{tabular}
\end{table}

The spatial coverage function for rangers at location $i$ follows a distance-decay model:

\begin{equation}
C_{spatial}^{ranger}(i,t) = \exp\left(-\frac{d_{i,station}}{\lambda}\right) \cdot \mathbb{I}(t \in \text{shift})
\label{eq:ranger_coverage}
\end{equation}

where $d_{i,station}$ is distance from nearest station and $\lambda \approx 8$ km is effective patrol radius.

For UAVs, spatial coverage incorporates flight patterns:

\begin{equation}
C_{spatial}^{UAV}(i,t) = \sum_{f \in \mathcal{F}} \mathbb{I}(i \in \text{flightpath}_f) \cdot \mathbb{I}(t \in \text{window}_f)
\label{eq:uav_coverage}
\end{equation}

\subsection{Threat Modeling}

Threat assessment incorporates historical poaching patterns, accessibility metrics, and species presence:

\begin{equation}
T_i = \sum_{s} \rho_s \cdot \text{Presence}_s(i) \cdot \left( \phi \cdot \text{Accessibility}_i + (1-\phi) \cdot \text{Historical}_i \right)
\label{eq:threat}
\end{equation}

where:
\begin{itemize}
    \item $\rho_s$ is species-specific vulnerability (higher for high-value targets)
    \item $\text{Accessibility}_i$ measures proximity to roads/boundaries (0-1 scale)
    \item $\text{Historical}_i$ is normalized past poaching incidence
    \item $\phi = 0.6$ weights accessibility vs. historical patterns
\end{itemize}

\subsection{Objective Function Decomposition}

The system-level protection objective decomposes into species-specific and area-weighted components:

\begin{equation}
\mathcal{P}_{system} = \omega_s \cdot \underbrace{\frac{\sum_{s} w_s \cdot P_s}{\sum_{s} w_s}}_{\text{Species protection}} + \omega_a \cdot \underbrace{\frac{\sum_{i \in \mathcal{A}} P_i}{|\mathcal{A}|}}_{\text{Area coverage}}
\label{eq:decomposition}
\end{equation}

with $\omega_s = 0.7$ and $\omega_a = 0.3$, reflecting emphasis on species conservation over pure area coverage.

This formulation enables direct comparison between uniform and strategic allocation strategies while maintaining biological relevance through the priority-weighted species protection component.
